In the present work, the following results were obtained:
\begin{itemize}
  \item It is shown that clouds in the Gromov--Hausdorff space are
    proper classes in the von Neumann--Bernays--G\"odel (NBG) axiom
    system, rather than sets.
  \item It is established that the Gromov--Hausdorff distance between
    clouds with a nontrivial intersection of stabilizer groups can
    take only the values \(0\) or \(\infty\).
  \item The theorem on the image of the center is proved: under a
    correspondence with finite distortion between the cloud of
    one-point spaces and a cloud with a nontrivial stabilizer, the
    image of the center lies at a distance not exceeding
    \(2\varepsilon\) (and hence finite).
  \item The cloud of the real line is investigated: the distances to
    the integer lattice and to other subspaces are computed.
  \item The notions of the angle and the isosceles angle of a cloud
    are introduced as new invariants of cloud geometry.
  \item It is proved that if clouds have a nontrivial intersection of
    stabilizers and distinct angles, then the distance between
    them is infinite.
\end{itemize}

In this work, we studied the Gromov--Hausdorff distance between
clouds --- equivalence classes of metric spaces with respect to
finiteness of the distance. We examined the properties of clouds from
the viewpoint of set theory and showed that they are proper classes
in the von Neumann--Bernays--G\"odel axiom system.

The structure of clouds with nontrivial stabilizer groups was
investigated. It was established that the distance between clouds
whose stabilizer groups have a nontrivial intersection is discrete
and takes only the values \(0\) or \(\infty\).

A new invariant --- the angle of a cloud --- was introduced and
studied, generalizing the geometric properties of the space of
bounded metric spaces. The values of the angle were computed for the
cloud of a one-point space and for the cloud of the real line. A
theorem relating the angles of clouds to the distance between them
was proved: it is shown that a difference in angles, together with a
nontrivial intersection of stabilizer groups, implies that the
Gromov--Hausdorff distance is infinite.

The obtained results refine the structure of the space of all metric
spaces from the perspective of the Gromov--Hausdorff metric on
unbounded objects. Limitations on the possibility of finite distance
between clouds of different geometric nature are revealed, and new
ways of classifying clouds via their angular invariants are demonstrated.
